\documentclass[12pt,a4paper]{article}

% --- 宏包加载 ---
\usepackage[UTF8]{ctex} % 中文支持
\usepackage{amsmath}    % 数学公式
\usepackage{amsfonts}
\usepackage{amssymb}
\usepackage{graphicx}   % 图片支持
\usepackage{geometry}   % 页边距设置
\usepackage{booktabs}   % 表格优化
\usepackage{circuitikz} % 逻辑电路绘图（可选，用于辅助说明）
\usepackage{float}      % 固定图片位置
\usepackage{caption}    % 图表标题

% 页边距设置
\geometry{left=2.5cm,right=2.5cm,top=2.5cm,bottom=2.5cm}

% --- 文档开始 ---
\begin{document}

\title{\textbf{数字电路实验报告：基于 74LS00 的三人表决器设计}}
\author{姓名：XXX \quad 学号：XXXXXXXX}
\date{\today}
\maketitle

\section{实验目的}
\begin{itemize}
    \item 掌握 74LS00（四 2 输入与非门）的基本逻辑功能。
    \item 学习组合逻辑电路的设计流程：真值表 $\rightarrow$ 逻辑表达式 $\rightarrow$ 电路图。
    \item 利用 Multisim 仿真软件进行功能验证及波形时序分析。
    \item 理解数字电路中的时延与毛刺现象，掌握异常分析与处理方法。
\end{itemize}

\section{实验原理}

\subsection{逻辑需求描述}
三人表决器共有三个输入端 $A$、$B$、$C$，代表三位投票人的意见（1 为同意，0 为反对）。根据“少数服从多数”原则，当有两个或两个以上的人同意时，输出 $F=1$；否则 $F=0$。

\subsection{真值表与逻辑推导}
根据需求列出真值表如下：

\begin{table}[H]
\centering
\caption{三人表决器真值表}
\begin{tabular}{ccc|c}
\toprule
$A$ & $B$ & $C$ & $F$ (输出) \\
\midrule
0 & 0 & 0 & 0 \\
0 & 0 & 1 & 0 \\
0 & 1 & 0 & 0 \\
0 & 1 & 1 & 1 \\
1 & 0 & 0 & 0 \\
1 & 0 & 1 & 1 \\
1 & 1 & 0 & 1 \\
1 & 1 & 1 & 1 \\
\bottomrule
\end{tabular}
\end{table}

由真值表写出主析取范式并简化：
$$F = AB\bar{C} + A\bar{B}C + \bar{A}BC + ABC$$
简化后的最简“与-或”表达式为：
$$F = AB + BC + AC$$

由于要求使用74LS00实现，根据摩根定律转换表达式：
$$F = \overline{\overline{AB + BC + AC}} = \overline{\overline{AB} \cdot \overline{BC} \cdot \overline{AC}}$$

\subsection{元器件说明}
74LS00 芯片内部包含四个 2 输入与非门。由于上述表达式包含三个“与非”项和一个三端“与非”项，实际接线时需利用 2 输入与非门进行级联。

\section{实验步骤与电路设计}

\subsection{电路拓扑设计}
1. 输入级：使用三个与非门分别计算 $\overline{AB}$、$\overline{BC}$ 和 $\overline{AC}$。
2. 输出级：将上述三个结果进行逻辑合成。由于 74LS00 仅有 2 输入端，三输入的“与非”操作实现为：
   $$\overline{X \cdot Y \cdot Z} = \overline{(X \cdot Y) \cdot Z} \text{ 需额外反相补偿}$$
   或采用如下变换：
   $$F = \overline{\overline{AB} \cdot \overline{BC} \cdot \overline{AC}} = \overline{\overline{AB \cdot BC} \cdot \overline{AC}} \text{ 利用两片 74LS00 共 6 个门完成}$$

\subsection{Multisim 仿真接线}
\begin{enumerate}
    \item 在 Multisim 中添加两个 \textbf{74LS00N} 芯片。
    \item 添加三个 \textbf{Interactive Digital Constant} 作为逻辑输入开关 A、B、C。
    \item 添加一个 \textbf{PROBE\_DIG\_BLUE} LED 灯作为输出指示。
    \item 按照推导的逻辑表达式完成连线，注意 VCC 和 GND 的正确接入。
\end{enumerate}

\section{实验分析与波形结果}

\subsection{功能验证分析}
通过逐一改变开关 A、B、C 的状态，观察 LED 灯的亮灭情况。仿真结果表明，当输入组合为 (011), (101), (110), (111) 时，LED 点亮，逻辑功能完全符合真值表要求。

\subsection{时域分析（Transient Analysis）}
设置时钟发生器输入 A、B、C 分别为不同频率的方波（例如 $f_A=100Hz, f_B=200Hz, f_C=400Hz$），进行瞬态仿真。
\begin{itemize}
    \item \textbf{预期波形}：在 $A, B$ 同时为高电平时，$F$ 立即翻转。
    \item \textbf{波形记录}：（此处应插入 Multisim 示波器截图）
    \item \textbf{时延说明}：由于 TTL 电路的平均传输延迟约为 $10ns$，在高速切换时，输出 $F$ 相对于输入的上升沿会存在微小滞后。
\end{itemize}

\subsection{频域与交流分析}
虽然表决器属于组合逻辑，但进行交流分析（AC Sweep）可以观察电路在极高频率下的增益稳定性。对于 74LS00，当工作频率超过 $30MHz$ 时，由于寄生电容影响，逻辑电平的上升/下降时间将占据信号周期的大部分，导致逻辑功能失效。

\section{异常记录与解决方案}

\begin{table}[H]
\centering
\caption{故障诊断表}
\begin{tabular}{p{4cm}|p{5cm}|p{5cm}}
\toprule
\textbf{异常现象} & \textbf{可能原因} & \textbf{解决方案} \\
\midrule
输出 LED 始终不亮 & 74LS00 的 VCC/GND 未接或芯片损坏 & 检查仿真模型电源连接，更换芯片 \\
\hline
输出逻辑在高频下紊乱 & 信号传输延迟（Race Hazard）产生毛刺 & 在输出端并联小电容滤波或优化布线 \\
\hline
输入悬空时输出异常 & TTL 门输入悬空默认为高电平 & 必须接 10k$\Omega$ 下拉电阻确保电平确定 \\
\bottomrule
\end{tabular}
\end{table}

\section{总结}
本实验成功利用 74LS00 设计并实现了三人表决器。通过逻辑代数变换，将“与-或”结构转为“与非-与非”结构。在 Multisim 仿真中，通过时域波形分析验证了逻辑转换的准确性。实验证明，数字电路设计中不仅要考虑逻辑正确性，还需关注硬件本身的电气特性（如延迟与悬空处理），这对于实际电路搭建具有重要的指导意义。

\end{document}